\subsubsection{\texttt{fnder}}

\texttt{fnder}计算一个样条函数的导函数，并仍以样条函数的形式返回。

\begin{itemize}
\item \texttt{ds = fnder(s, order)} 得到样条函数 \texttt{s} 的 \texttt{order} 阶导函数，\texttt{s} 可以是 pp 格式、B 格式。
\item \texttt{order} 为正数代表对原函数求导数，\texttt{order} 为负数代表对原函数求积分（积分常数设为0）。
\item \texttt{order} 是一个不超过样条函数次数的正整数。
\item 该函数输出的形式与输入相同，仍以样条函数的格式返回，结果要么是 \texttt{pp} 格式，要么是\texttt{B} 格式。
\item \texttt{fnder(s)}同\texttt{fnder(s, 1)}。
\end{itemize}

如输入：
\begin{lstlisting}[language = matlab]
  x  = [   0, 2, 3,  5, 6  ]; 
  y  = [1, 0, 3, 1, -1, 0, 0]; 
  s  = csape(x,y,'complete');
  ds = fnder(s); 
\end{lstlisting}

输出为：
\begin{lstlisting}[language = matlab ]
  form: 'pp'
  breaks: [0 2 3 5 6]
  coefs: [4 * 3 double]
  pieces: 4
  order: 3
  dim: 1
\end{lstlisting}
\textbf{B格式下\texttt{fnder}函数实现说明:}
B格式下\texttt{fnder}函数实现分为两个部分：求导和求积分。

记输入的节点序列为$[t_0,t_1,\dots,t_m,t_{m+1}]$的B格式样条函数为
\begin{equation}
  \label{eq:given-b-spline-fnder}
  S(x) = \sum_{i=1}^{m-n+1} c_i B_i^{n}(x)
\end{equation}
根据B样条基函数的求导规则（\textcolor{blue}{这里要求$n \geq 2$，暂时如此实现}）
\begin{equation}
  \label{eq:derivatives-of-a-basis}
  \difFrac{}{x} B_i^n(x) = \frac{n}{t_{i+n-1} -t_{i-1}} B_i^{n-1}(x) - \frac{n}{t_{i+n} -t_i}B_{i+1}^{n-1}(x)
  := \alpha_i B_i^{n-1}(x) - \beta_i B_{i+1}^{n-1}(x)
\end{equation}
对$S(x)$进行求导可得
\begin{equation}
  \begin{aligned}
    S^\prime(x) &= \sum_{i=1}^{m-n+1} c_i \left[B_i^{n}(x)\right]^\prime \\
                &= \sum_{i=1}^{m-n+1} c_i \left[\alpha_i B_i^{n-1}(x) - \beta_i B_{i+1}^{n-1}(x)\right] \\
                &= c_1 \alpha_1 B_1^{n-1}(x) + \sum_{i=2}^{m-n+1}(c_i \alpha_i -c_{i-1}\beta_{i-1})B_{i}^{n-1}(x) 
                  - c_{m-n+1} \beta_{m-n+1} B_{m-n+2}^{n-1}(x)
  \end{aligned}
\end{equation}
根据以上推导可以实现B格式下\texttt{fnder}函数求导数的部分。\textcolor{blue}{\kaishu（需要注意的一点是，因为重节点的存在，以上
  公式中$B_{i}^{n-1}(x),B_{m-n+2}^{n-1}(x)$的结点可能是全部相同的，根据基函数的定义，这样的B样条基函数就是0函数，MATLAB的
  做法是在$S^\prime(x)$的计算结果中将这样的基函数剔除，于是会出现求导过后反而构成$S(x)$的基函数个数变少的情况。）}

对于求积分问题，设
\[
  S(x) = \sum_{i=1}^{\text{Number}} y_i B_i^{n},
\]
通过在右侧添加重节点的方式，将 $S$ 的基函数进行延拓，求 $S$ 的原函数等价于求
\[
  S + 0 * \sum_{i=1}^{\text{Order + 1 - multinum}} B_{\text{Number + i}}^{n}
\]
的原函数。其中 multinum 为最右侧重节点数，当最后一个点非重节点时 multinum $=1$，若重节点过多，则报错。

接下来考虑 $S$的原函数，在上述得到的节点最后再添加一个重节点。则可以以
\[
  \left\{ B_i^{n+1} \right\}_{1}^{\text{Number + Order + 1 - multinum}}
\]
为基函数进行待定系数，再根据B-spline求导规则 式 (\ref{eq:derivatives-of-a-basis})，可以列出线性方程组。关键在于此时
\[
  \difFrac{}{x} B_{last}^{n+1} = \frac{n+1}{t_{last + n + 1 - 1 } - t_{last -1}} B_{last}^{n} + 0.
\]
该方程组为：
\[
  \begin{pmatrix} 
    \alpha_{1}   &  &  &  \\
    -\alpha_{2}  & \alpha_{2} &  &  \\
             & -\alpha_{3} & \ddots &  \\
             &  & \ddots & \alpha_{N-1} \\
             &  &  & \alpha_{N}
  \end{pmatrix}
  \begin{pmatrix} 
    x_{1} \\
    x_{2} \\
    \vdots \\
    x_{N-1} \\
    x_{N}
  \end{pmatrix}
  =
  \begin{pmatrix} 
    y_{1} \\
    y_{2} \\
    \vdots \\
    y_{N-1} \\
    y_{N}
  \end{pmatrix}.
\]

